\documentclass[12pt]{article}
\usepackage[margin=0.6in,top=0.85in,bottom=0.55in]{geometry}
\usepackage{lmodern}
\usepackage{microtype}
\usepackage{fancyhdr}
\usepackage{titlesec}
\usepackage{enumitem}
\usepackage{lastpage}
\usepackage[hidelinks]{hyperref}

\setlength{\parindent}{0pt}
\setlength{\parskip}{0.25em}
\setlength{\headheight}{26pt}
\titleformat{\section}{\large\bfseries\scshape}{}{0pt}{}
\pagestyle{fancy}
\fancyhf{}
\fancyhead[C]{\footnotesize Student Name: \hspace{2.3in} \quad Assignment ID: U1L8LE \quad Page \thepage\ of \pageref{LastPage}}
\renewcommand{\headrulewidth}{0.5pt}

\newcommand{\answerbox}[2]{%
\vspace{0.12em}
\noindent\textbf{#1}\par
\vspace{0.04em}
\noindent\fbox{%
\begin{minipage}[t][#2]{\dimexpr\textwidth-2\fboxsep-2\fboxrule}
\rule{\textwidth}{0pt}
\vspace{#2}
\end{minipage}}
\par\vspace{0.22em}}

\begin{document}

\begin{center}
{\LARGE\bfseries Week 8 Lit Example Worksheet}\par\vspace{0.05em}
{\large\scshape Macbeth: Fair, Foul, and Shifting Security}\par\vspace{0.12em}
\rule{0.78\textwidth}{0.5pt}
\end{center}

\textbf{Purpose:} Apply Week 8's method to \textit{Macbeth}. Expose shifting terms and retest the inference with fixed meanings.

\textbf{Directions:} Use the Lit Example Reader. Do not retell the whole plot.

\section*{Part 1: Fair and the Greetings}

\textbf{Part 1 Q1:} In Passage 1, give two ordinary senses of \textit{fair} that the language can carry.

\answerbox{Part 1 Q1 ANSWER BOX}{2.15in}

\textbf{Part 1 Q2:} Which greetings are present facts, and which are future claims? Why does that difference matter for a stable argument?

\answerbox{Part 1 Q2 ANSWER BOX}{2.35in}

\newpage

\section*{Part 2: From Sounding Fair to Acting}

\textbf{Part 2 Q1:} Reconstruct the risky inference Macbeth may draw from the witches' words in a middle-term shape (even if informal). Name the term that is doing the linking work.

\answerbox{Part 2 Q1 ANSWER BOX}{2.45in}

\textbf{Part 2 Q2:} Explain why that link is an ambiguous middle or term shift rather than a stable proof that he must be king by crime.

\answerbox{Part 2 Q2 ANSWER BOX}{2.35in}

\newpage

\section*{Part 3: Security Promises}

\textbf{Part 3 Q1:} For ``none of woman born shall harm Macbeth,'' state Sense A (how Macbeth hears it) and Sense B (a different sense that later undercuts the guarantee).

\answerbox{Part 3 Q1 ANSWER BOX}{2.45in}

\textbf{Part 3 Q2:} For the Birnam wood promise, state Sense A and Sense B, then say whether Macbeth's ``That will never be'' still follows under both fixed senses.

\answerbox{Part 3 Q2 ANSWER BOX}{2.45in}

\newpage

\section*{Part 4: Repair}

\textbf{Part 4 Q1:} Rewrite one of the security arguments with fixed meanings so that it no longer falsely guarantees Macbeth's safety.

\answerbox{Part 4 Q1 ANSWER BOX}{2.55in}

\textbf{Part 4 Q2:} In one paragraph, explain how Week 8's tool changes a reading of these passages from ``magic prophecy'' to ``unstable terms used as if they were proof.''

\answerbox{Part 4 Q2 ANSWER BOX}{2.45in}

\end{document}
